\documentclass[10pt,twocolumn]{report}
\usepackage[margin=0.75in]{geometry}
\usepackage{amsmath,amssymb,amsthm}
\usepackage{booktabs}
\usepackage{array}
\usepackage{graphicx}
\usepackage{adjustbox}
\usepackage{lmodern}
\usepackage{microtype}
\usepackage{seqsplit}
\usepackage{hyperref}
\usepackage[numbers,sort&compress]{natbib}

% Two-column readability: let long Lean identifiers break at underscores,
% scoped to \code only so bib keys / URLs / bookmarks are untouched.
\newcommand{\code}[1]{\texttt{\seqsplit{#1}}}
\AtBeginDocument{%
  \renewcommand{\[}{\par\smallskip\noindent\begin{adjustbox}{max width=\columnwidth}$\displaystyle}%
  \renewcommand{\]}{$\end{adjustbox}\par\smallskip}%
}
\emergencystretch=3em
\sloppy
\newcommand{\Ste}[1]{\texttt{\seqsplit{Ste.#1}}}

\newtheorem{theorem}{Theorem}[chapter]
\newtheorem{proposition}[theorem]{Proposition}
\newtheorem{corollary}[theorem]{Corollary}
\theoremstyle{definition}
\newtheorem{definition}[theorem]{Definition}
\theoremstyle{remark}
\newtheorem{remark}[theorem]{Remark}
\newtheorem{interpretation}[theorem]{Interpretation}

\title{Set-Theoretic Estimation over Non-Metric Domains\\
\large From Feasibility in Hilbert Space to Coherence in Language:\\
a Machine-Checked Monograph}
\author{Project \texttt{thejeb}}
\date{\today}

\begin{document}
\maketitle

\begin{abstract}
Set-theoretic estimation (STE) answers an estimation problem by
intersection: an estimate is acceptable exactly when it is consistent
with every piece of available information.  The classical theory
\citep{combettes1993} lives in metric --- indeed Hilbert --- space,
where convexity and projection operators make the feasible set both
meaningful and computable.  This monograph is about what happens to STE
\emph{when the underlying set is non-metric}.  We follow
\citet{carlson2012} in the topological origination of STE --- set
estimation over topological spaces with no metric, motivated by
decryption --- and then carry the theory into a domain where no metric
was ever on offer: language.  The carriage proceeds in three moves:
first, a link to \emph{parsing}, where Constraint Grammar's eliminative
algebra is shown to be exactly STE; second, \emph{frame analysis},
where set-valued and dynamic frames give a possible-worlds semantics
for machine-gathered meaning; third, accepting the challenge of
\emph{multi-document coreference}, where identity itself becomes a
lattice limit --- the least partition adequate to the accumulated
evidence.  What replaces the metric is order: the feasibility lattice,
its sheaf-theoretic gluing, and a combinatorial Shannon geometry on
readings.  The monograph's unifying thread is a single machine-checked
dichotomy in that geometry: joining information is a metric contraction,
extracting common information is not --- and this join/meet conditioning
asymmetry recurs, under three other names, as the tractable/hard,
decomposed/coupled, and gluable/obstructed divides that organize every
chapter.  Every result cited as a theorem is verified in the project's
Lean~4 library and cited by exact identifier; everything else is
framing, and is marked as such.
\end{abstract}

\tableofcontents

\chapter*{Preface: how to read this book}
\addcontentsline{toc}{chapter}{Preface}

This monograph is a \emph{spine}, not an omnibus.  The project it
unifies has produced a body of papers and notes, each of which remains
the owning, self-contained deep dive for its topic; this book supplies
the connective narrative that makes them chapters of one argument, and
it originates in full only the material that exists nowhere else in
prose --- chiefly the multi-document coreference theory of
Chapter~\ref{ch:coref}.  Each part opens with synthesis that is new
here, states what its topic contributes to the whole, and then
delegates: ``the full treatment is [paper].''  Appendix~\ref{app:map}
gives the complete ownership map, together with recommendations for
trimming the residual overlaps among the existing papers.

Three disciplines hold throughout.

\emph{First, the mechanization bright line.}  Every claim labeled a
theorem is machine-checked in the project's Lean~4 / Mathlib
\citep{moura2021,mathlib2020} library, sorry-free, verified in
continuous integration, and cited by its exact Lean identifier, e.g.\
\code{frame\_hasHelly\_two} (\Ste{HellyConsistency}).  Surrounding
prose is interpretation --- ours to defend, never the kernel's.

\emph{Second, coherence, not truth.}  The theory studies when a body of
machine-gathered readings is \emph{coherent} --- jointly satisfiable ---
not whether it is \emph{true}.  The epistemological position, its
bracketing of rival pictures of meaning, the fidelity/consistency
conditional (a feasible set is only as good as the extraction that
produced its constraints), and the Shannon-entropy ceiling on what
crisp set-valued methods can express are stated once, carefully, in the
project's epistemology note (\emph{Hyperframes and the Epistemology of
Machine-Gathered Semantics}, \code{hyperframes-epistemology.tex}).  We
reference that position throughout and do not re-litigate it.

\emph{Third, no duplicated deep dives.}  Where a chapter's material has
an owning paper, this book states the result, its place in the arc, and
its Lean anchor --- and stops.

\part{From Metric to Topological: the Classical Inheritance}
\label{part:classical}

\chapter{Estimation by intersection: the metric theory}
\label{ch:combettes}

Set-theoretic estimation begins with a reversal of habit.  Classical
estimation optimizes: it posits a criterion and returns the single
point that extremizes it.  STE instead \emph{intersects}: every piece
of information --- a measurement, a bound, a prior structural fact ---
is recorded as the set of hypotheses consistent with it (its
\emph{property set}), and the deliverable is the \emph{feasibility
set}, the intersection of all property sets.  Any member is an
acceptable estimate; the set itself is the honest answer.
\citet{combettes1993} founded the framework in this generality, drawing
on a lineage running through unknown-but-bounded state estimation
\citep{schweppe1968}, image recovery by convex projections
\citep{youla1982,combettes1993numerical}, and set-membership
identification in signal processing \citep{deller1989}.

The classical theory is emphatically \emph{metric}.  Its home is a
Hilbert space; its property sets are closed and convex; its engine is
the projection operator, well-defined precisely because a nearest point
exists and is unique in that geometry.  The great algorithmic corpus
--- POCS, parallel projection methods, and their convergence theory
\citep{bauschke1996,combettes2011proximal} --- is a corpus about
metric structure: distances to sets decrease, Fej\'er monotonicity
holds, fixed points attract.  When textbooks say ``set-theoretic
estimation works,'' they mostly mean: \emph{the metric makes it work}.

This monograph asks what survives when the metric is removed --- and
the honest first answer is: the \emph{definitions} survive untouched.
Feasibility is an intersection; intersections need no metric.  The
consistent/fair/ideal taxonomy of estimators, information monotonicity
(more constraints, smaller feasible set), and the detection of invalid
information by inconsistency are order-theoretic facts.  The project's
machine-checked core formalizes exactly this metric-free skeleton ---
the feasibility set and its lattice laws, the taxonomy, antitone
refinement --- and verifies that the classical statements hold with no
appeal to distance.  \textbf{The full treatment is the mechanization
paper} (\emph{A Machine-Checked Core for Set Theoretic Estimation},
\code{ste-mechanization.tex}), which owns the formal model and the
core theorem list; the primary-source assessment of
\citet{combettes1993} is owned by the literature review
(\code{literature-review.tex}).

What does \emph{not} survive metric removal is the algorithmics: no
projections, no nearest points, no convergence-by-contraction.  The
burden of the rest of this book is that something else takes their
place --- order, lattice, and sheaf structure --- and that the
replacement is not a consolation prize but a theory with its own exact
dichotomies, its own geometry, and its own frontier.

\chapter{The topological origination: Carlson}
\label{ch:carlson}

The first deliberate step off the metric was taken by
\citet{carlson2012}, who recast STE over bare topological spaces to
attack decryption.  The move is conceptually prior to any application:
if property sets are closed sets of a compact space, the feasibility
set is closed, and nonemptiness --- existence of an estimate --- follows
from the finite-intersection property, with no metric anywhere.
Carlson's cryptanalytic instantiation makes the stakes concrete: the
hypothesis space is a key space, each intercepted ciphertext constrains
it, and estimation is the shrinking of the feasible key set; the
unicity phenomenon of \citet{shannon1949} becomes a statement about
when the feasibility set collapses to a point.  A continuing line of
work develops the program
\citep{carlson2005,carlson2014breaking,carlson2021keyspace,
carlson2022equivalence,carlson2022localentropy,carlson2026entropy}.

For this monograph Carlson supplies three things.  First, the
\emph{license}: STE over non-metric spaces is not our invention but an
established origination that we follow and mechanize --- the project's
library formalizes the closedness/compactness existence theorems, the
skeleton of Carlson's STE-decryption reduction, and his key-counting
and cipher-reduction chain (owned by \code{ste-mechanization.tex}).
Second, the \emph{partition insight}: Carlson's asymptotic-equipartition
work partition is STE-invariant --- an early sign that when the metric
goes, \emph{partitions} (and their lattice) become the load-bearing
structure; this observation is developed in the language note
(\code{ste-and-language.tex}) via rough-set indiscernibility
\citep{pawlak1982rough}.  Third, the \emph{template for a domain
carriage}: Carlson carried STE from signal recovery to cryptanalysis by
identifying the domain's native constraint structure.  Parts
\ref{part:parsing}--\ref{part:frontier} of this book perform the same
carriage into language.

The surveys own the surrounding scholarship: the projection-algorithm
inheritance and the set-membership tradition
(\code{topological-ste-survey.tex}), and the post-2012 uptake of
topological STE together with the mechanization gap that this project
fills (\code{decade-followon-survey.tex}).  \textbf{Those two surveys,
plus the literature review, are the owning deep dives for everything
historical in this part}; we do not repeat their bibliographies here.

One more inheritance deserves its own flag, because it becomes the
spine of the book.  Shannon himself, in a paper that predates STE by
four decades, described \emph{the lattice theory of information}
\citep{shannon1953lattice}: information sources ordered by refinement,
with join and meet.  Modern work \citep{delsol2024} equips that lattice
with a metric derived from conditional entropy.  When STE goes
non-metric, it lands --- provably, in the finite counting shadow --- on
exactly Shannon's lattice.  That landing is Chapter~\ref{ch:spine}.

\part{The Non-Metric Structure Theory}
\label{part:structure}

\chapter{The conditioning dichotomy: the spine of the book}
\label{ch:spine}

This chapter states, once, the single mathematical fact that organizes
the entire monograph.  Later chapters will not restate it; they will
point here.

\section{The lattice that replaces the metric space}

Remove the metric from STE and ask what structure the readings of a
corpus actually carry.  The project's answer, accumulated across its
Lean modules, is: a \emph{partition lattice with a counting geometry}.
Hypotheses about a corpus induce partitions (which claims corefer,
which readings agree); partitions form a complete lattice; and on that
lattice there is a genuine metric --- not assumed, but derived.  Define
the combinatorial Shannon distance as the counting shadow of
$D(X,Y)=H(X\mid Y)+H(Y\mid X)$:
\[
  \code{shannonDist}\,(P,Q)
  \;=\;
  \bigl(r(P \sqcup Q)-r(P)\bigr)+\bigl(r(P\sqcup Q)-r(Q)\bigr),
\]
where $r$ is the partition rank (\code{shannonDist},
\Ste{InfoDistance}; \code{rank}, \Ste{PartitionRank}).  The triangle
inequality is a theorem, and its proof is the submodularity of rank
(\code{shannonDist\_triangle}, \Ste{InfoDistance};
\code{rank\_submodular}, \Ste{PartitionRank}) --- the lattice-theoretic
residue of the entropy inequalities, exactly as in the
information-lattice synthesis of \citet{delsol2024} descending from
\citet{shannon1953lattice}.

The honest first observation kills a tempting misreading: on a finite
carrier this metric is uniformly discrete --- distinct partitions are at
distance at least one (\code{one\_le\_shannonDist\_of\_ne},
\Ste{InfoTopology}) --- so the induced point-set \emph{topology} is
trivial.  The content of the non-metric theory is not topological but
\emph{metric-geometric}: it lies in how the lattice operations behave
with respect to this distance.

\section{The dichotomy}

Here is the spine.  The two lattice operations are \emph{not}
symmetric with respect to the Shannon geometry.

\begin{theorem}[Join is a contraction]
For all partitions $P,Q,R$,
\[
  d(P\sqcup R,\; Q\sqcup R) \;\le\; d(P,Q),
\]
and join is jointly 1-Lipschitz:
$d(P\sqcup Q, P'\sqcup Q') \le d(P,P') + d(Q,Q')$.
\emph{(\code{shannonDist\_sup\_right\_le},
\code{shannonDist\_sup\_le\_add}, \Ste{InfoTopology}.)}
\end{theorem}

\begin{theorem}[Meet is not]
There exist partitions $P,Q,R$ of a four-element carrier with
$d(P,Q) < d(P\sqcap R,\; Q\sqcap R)$; on a $2k$-element carrier the gap
scales with $k$, so meet admits no Lipschitz constant uniform in the
carrier.  \emph{(\code{shannonDist\_inf\_not\_contraction},
\Ste{InfoTopology}; the scaling claim is interpretation of the
construction, the $k=2$ instance is the mechanized witness.)}
\end{theorem}

\begin{interpretation}
\emph{Aggregating} information --- join, ``all these readings
together'' --- is well-conditioned: perturbing the inputs moves the
aggregate by no more than the perturbation.  \emph{Extracting what
readings share} --- meet, common information --- is ill-conditioned: an
arbitrarily small change in the inputs can move the common part
arbitrarily far.  This is the counting shadow of the G\'acs--K\"orner
phenomenon --- common information is discontinuous and far less than
mutual information \citep{gacs1973common} --- surfacing inside STE with
no probability in sight.
\end{interpretation}

The remaining geometry sharpens the picture: the metric is
bipartite-parity like a hypercube skeleton
(\code{shannonDist\_add\_rank\_add\_rank}) yet embeds in no $\ell_1$
cube (\code{pentagonal\_violation}) and is not median
(\code{median\_not\_unique}); it is geodesic along refinement chains
(\code{aligned\_of\_chain}, \code{le\_sup\_of\_aligned}) --- the
covering-graph distance of a geometric lattice, genuinely its own
space (all \Ste{InfoTopology}).  Coherence-space cannot be
coordinatized away into a Hamming cube even in its crisp counting
shadow.

\section{The same dichotomy under three other names}

The claim that earns this chapter its title: the join/meet conditioning
asymmetry is the \emph{same} divide that the structure theory meets
three more times, each occurrence machine-checked.

\paragraph{Tractable versus hard.}  For frame-structured corpora,
pairwise consistency forces global consistency: the Helly number is two
(\code{frame\_hasHelly\_two}, \Ste{HellyConsistency}).  Checking
coherence is a \emph{join-side} activity --- accumulate constraints,
scan pairs --- and it is correspondingly cheap and stable.  The moment
a question requires knowing what a family of readings has \emph{in
common} at higher order, locality fails: for second-order (set-of-sets)
sources the Helly number provably jumps to at least three
(\code{exists\_not\_hasHellySel\_two}, \Ste{Carlson.SetOfSets};
Chapter~\ref{ch:frames}).

\paragraph{Decomposed versus coupled.}  The partition rank is
submodular (\code{rank\_submodular}), and the \emph{coupling excess}
of a constraint against a product structure is zero exactly when the
readings decompose and strictly positive exactly when they genuinely
entangle (\code{couplingExcess\_pos\_of\_not\_le},
\Ste{CouplingRank}).  Decomposable constraints are the ones whose
information content is exhausted by joins of marginal information;
coupling is precisely the residue that only a meet-like common-part
analysis would see --- and it is the exact source of representational
blow-up (Chapter~\ref{ch:representation}).

\paragraph{Gluable versus obstructed.}  Rectangular (product)
constraints glue over every cover with vanishing \v{C}ech obstruction
(\code{propertySet\_cechVanishesCover}, \Ste{SetValuedFrame}); the
diagonal --- the archetypal coupled constraint --- has obstruction
exactly~$2$ (\code{diagonal\_cechObstruction}, \Ste{CechObstruction}).
Local-to-global success is the geometric face of decomposition; the
obstruction is the geometric face of coupling
(Chapter~\ref{ch:gluing}).

\begin{interpretation}
One dichotomy, four costumes: join/meet conditioning
(well/ill-conditioned), Helly locality (tractable/hard), coupling
excess (decomposed/coupled), \v{C}ech obstruction
(gluable/obstructed).  In each case the benign side is the side
reachable by accumulating and combining --- the join side --- and the
malignant side is the side that asks what parts share, agree on, or
jointly determine --- the meet side.  Every subsequent chapter of this
book instantiates the dichotomy in its own domain, and none of them
escapes it.
\end{interpretation}

The detailed geometry --- proofs-in-prose, the non-embeddability
witnesses, the asymptotic Rokhlin-metric outlook --- is recorded in the
interpretive log (\code{hyperframe-interpretations.tex}, the
information-lattice section), which is the owning document for
\Ste{InfoTopology} until its dedicated paper exists; this chapter is
the monograph-level statement and is deliberately complete as such.

\chapter{Coherence as gluing: sheaves, obstructions, and the gluing complex}
\label{ch:gluing}

When the hypothesis space is a product over $N$ variables, the
feasibility set lives behind an exponential wall, and every practical
scheme keeps only \emph{local} information: small tables over small,
overlapping sets of variables.  Every such scheme lives or dies on one
question --- when does locally consistent partial information determine
a global solution?  The mathematics of that question is sheaf theory
and \v{C}ech cohomology, transplanted to constraint satisfaction by the
contextuality program of
\citet{abramsky2011sheaf,abramsky2012cohomology}, and developed for STE
in machine-checked form by this project's flagship paper.

\textbf{The full treatment is the flagship} (\emph{Local Pieces, Global
Solutions}, \code{ste-cohomology.tex}), whose centerpiece, the
support--cover gluing theorem (\Ste{SupportCover}), characterizes
exactly when every compatible family of local sections glues: precisely
when the constraint is generated by pieces whose scopes fit inside the
cover.  Its accessible companion (\code{ste-cohomology-recast.tex})
re-derives the dictionary for readers with no topology; the
transplant's relation to the contextuality literature is owned by the
obstruction lit-review (\code{cohomological-obstruction-litreview.tex}).
Within the spine's terms: gluing succeeds on the decomposed side of the
dichotomy and is obstructed on the coupled side, with the diagonal's
obstruction of exactly $2$ (\code{diagonal\_cechObstruction}) as the
minimal witness.

What this chapter adds --- because it exists so far only in Lean and in
the interpretive log --- is the \emph{refinement} of the gluing verdict
into a shape.  The module \Ste{LocalGluing} drops the all-or-nothing
sheaf condition: \code{GluesOn} asks whether a compatible family
assembles over a \emph{sub-cover}, and the family of sub-covers on
which it does --- the \code{gluingComplex} --- is downward closed
(\code{mem\_gluingComplex\_of\_subset}), hence an abstract simplicial
complex.  The tight sheaf condition is exactly the presence of the top
face (\code{cechVanishesCover\_iff\_gluesLocallyAt\_card}), and the
diagonal separates the notions: it glues locally at width one but not
at width two
(\code{diagonal\_gluesLocallyAt\_one},
\code{diagonal\_not\_gluesLocallyAt\_two},
\code{diagonal\_local\_not\_tight}).  A bridge theorem returns to the
spine's tractable side: enough Helly structure on the gluing fibers
forces the top face --- local coherence climbs to global
(\code{cechVanishesCover\_of\_hasHelly\_gluingFiber}).

\begin{interpretation}
Incoherence is not a single bit.  A corpus that fails to cohere
globally still coheres over a structured family of sub-corpora, and
that family --- the gluing complex --- is the real object of study: the
paraconsistent stance \citep{dacosta1974paraconsistent} made geometric.
The obstruction is thereby \emph{located}, not merely detected.  This
becomes load-bearing for language in Part~\ref{part:frames}: a corpus
of documents that contradict each other still has a shape of partial
coherence, and Chapter~\ref{ch:empirical}'s ratification corpus
exhibits exactly that shape in the wild.
\end{interpretation}

The gluing complex is new material; its deep dive belongs to a future
revision of the flagship (Appendix~\ref{app:map} records this
assignment), and this section is its monograph-level statement.

\chapter{Representation and its limits}
\label{ch:representation}

The gluing theory says \emph{when} local pieces determine the whole;
the representation theory says \emph{what it costs}.  The project's
answer is a mechanized dichotomy stack: rectangular constraints are the
tractable upper bound; coupling is the obstruction
(\code{couplingExcess\_pos\_of\_not\_le}); block decomposition trades
sums for products; conditioning and variable elimination are
machine-checked; treewidth bounds the bags; and unbounded coupling is
provably exponential --- with the relation between obstruction size and
representation blow-up a \emph{sandwich, not an equality}.
\textbf{The full treatment is the representation paper}
(\emph{Representation and Tractability in Set Theoretic Estimation},
\code{ste-representation-sheaves.tex}), with the sandwich's
literature setting owned by \code{representation-bounds-litreview.tex};
the classical constraint-processing background is
\citet{dechter2003constraint,dechter1999bucket,robertson1986graphminorsII,
yannakakis1991expressing}, and the lower-bound toolkit descends from
communication complexity and extension complexity
\citep{kushilevitz1997communication,fiorini2015exponential}.

In the spine's terms this chapter is the computational shadow of the
dichotomy: everything join-like (accumulate, eliminate, decompose along
a tree) is polynomial in the width; everything meet-like (what do these
variables jointly determine that no projection sees) is where the
exponential lives.  We state this once and move on; the paper owns the
proofs, and Appendix~\ref{app:map} flags its one overlap with the
flagship for trimming.

\part{Language I: Parsing}
\label{part:parsing}

\chapter{Constraint Grammar is set-theoretic estimation}
\label{ch:parsing}

The first move into language is the most conservative one available:
find the corner of computational linguistics that is \emph{already}
STE, and prove it.  That corner is Constraint Grammar
\citep{karlsson1990constraint,karlsson1995constraint}: begin with a
universe of candidate morphosyntactic analyses, and let rules
\emph{eliminate} candidates until what survives is the parse.  The
practice matured through implemented parsers
\citep{tapanainen1994tagging,bick2015cg3}, always described
operationally, never estimation-theoretically.

The identification is now a theorem-level equivalence: for the hard,
purely eliminative fragment, a CG rule \emph{is} an STE property set,
the surviving analyses \emph{are} the feasibility set, rule addition is
antitone refinement, fairness is preservation of the intended parse,
and unambiguous successful analysis is ideal information --- verified
as an explicit equivalence in \Ste{ConstraintGrammar}.  \textbf{The
full treatment is the parsing note} (\emph{The Algebraic Core of
Constraint Grammar as Set Theoretic Estimation},
\code{constraint-grammar-ste.tex}), which also owns the scoping: CG
operations that genuinely change the candidate universe are not
isomorphic to static STE and must be treated as staged maps between STE
problems.  The broader historical claim --- that unification grammar
\citep{shieber1986unification} and underspecification semantics
\citep{copestake2005mrs,egg2001cls,bos1996hole} are likewise doing
uncited STE, and that the bridge to \citet{combettes1993} is unclaimed
--- is owned by the language note (\code{ste-and-language.tex}),
together with the Formal Concept Analysis reading
\citep{ganter1999fca} in which feasible sets are extents of a concept
lattice.

For the monograph's arc, parsing contributes two lessons.  First,
\emph{proof of carriage}: language admits STE not by analogy but by
isomorphism, in at least one exactly-delimited fragment --- the
non-metric theory has a native application, with the domain supplying
eliminative constraints where cryptanalysis supplied ciphertexts.
Second, \emph{the boundary lesson}: the isomorphism stops exactly where
the candidate universe stops being fixed.  A parser's universe is one
sentence's analyses; a corpus's universe of \emph{readings} is neither
fixed nor crisp.  Pushing past that boundary is what forces the
set-valued frame --- the primitive of Part~\ref{part:frames}.

\part{Language II: Frames}
\label{part:frames}

\chapter{The set-valued frame and its deformations}
\label{ch:frames}

The primitive of the language-side theory is not the crisp assertion
``$X=v$'' but the \emph{possibility set}: a source commits only to
``the referent lies in $S$.''  The epistemology note
(\code{hyperframes-epistemology.tex}) owns the defense of this choice
--- reference as gathered by machines or readers is underdetermined but
bounded, and pretending otherwise is false to the gathering --- along
with everything the position brackets.  Here we record what the
primitive \emph{is}, machine-checked, and how it deforms.

The set-valued frame is a first-class object (\Ste{SetValuedFrame}):
each variable receives a possibility set; the induced constraint is the
product box, hence rectangular (\code{propertySet\_isRectangle}) and
obstruction-free over every cover
(\code{propertySet\_cechVanishesCover}); a corpus is feasible exactly
when every variable's possibility-set intersection is nonempty
(\code{corpusFeasibilitySet\_nonempty\_iff}); the feasibility set is a
formal-concept extent (\code{isExtent\_corpusFeasibilitySet}),
unifying the possibility-set and concept-lattice
\citep{ganter1999fca} pictures; and the crisp determinate-reference
model embeds as the singleton case with 2-locality transported
(\code{hasHelly\_ofCrisp\_two}).  In spine terms: the set-valued frame
is engineered to live entirely on the benign side of the dichotomy ---
rectangular, gluable, Helly-2 --- which is precisely why frame-based
corpus checking is tractable at all.

Two deformations test how far the benign side extends.

\emph{Graded possibility.}  A fuzzy frame \citep{zadeh1965fuzzy} grades
membership in $[0,1]$; feasibility becomes a
\code{consistencyDegree}; and the $\alpha$-cut theorem
(\code{cut\_fuzzyFeasibility}, \Ste{FuzzyFrame}) shows the graded
theory is a \emph{stack} of crisp theories indexed by confidence
threshold --- every crisp result acquires a graded shadow by cutting,
the crisp theory sits at the boundary
(\code{consistencyDegree\_crisp\_eq\_one\_iff}), and 2-locality
survives in graded form (\code{weightedFrame\_gradedHelly\_two}).
Grading is a \emph{safe} deformation.

\emph{Higher-order possibility.}  A Carlson-style set-of-sets source
--- ``the answer is one of these sets'' --- makes consistency a
selectability problem (\code{SelectableConsistent},
\Ste{Carlson.SetOfSets}).  The singleton case recovers the first-order
theory exactly (\code{selectable\_singleton\_iff},
\code{hasHellySel\_singletonFamily\_iff}), but the deformation is
\emph{not} safe: a three-element triangle of two-element families is
pairwise selectable with no global selection
(\code{exists\_not\_hasHellySel\_two}), so the second-order Helly
number is at least three.  Disjunction over readings costs locality ---
the first recurrence of the spine's dichotomy on language soil, and an
operational correctness cliff: any solver ingesting set-of-sets data
that relies on a pairwise scan is unsound.

\textbf{The owning document for both deformations and their
interpretation is the interpretive log}
(\code{hyperframe-interpretations.tex}), pending their dedicated
paper; the epistemology note owns the position that the probabilistic
deformation (a normalized density with Shannon entropy on top) is
currently out of mechanized reach, and why the honest boundary is drawn
there.

\chapter{Dynamic corpora: frames in motion}
\label{ch:dynamic}

Real corpora are mutable: documents arrive, are deleted, and carry
several candidate frame interpretations at once; coreference among
their claims must not be forced to one early clustering.  The project's
answer is a possible-worlds semantics: every exact interpretation and
strict coreference partition is a candidate normalization hypothesis,
constraints denote property sets of hypotheses, and the solver is
literally STE.  Identity statements then come in exact modal grades ---
inconsistent, must, cannot, uncertain --- with must-coreference the
safe quotient and the uncertainty overlay provably non-transitive.
\textbf{The full treatment is the dynamic-frames paper}
(\emph{Dynamic Set-Theoretic Normalization}, %
\code{dynamic-frame-normalization.tex}): document-insertion
antitonicity, deletion expansion, order independence, delta
decomposition, minimal document supports, the $2^N$ tight limit, and an
executable solver proved extensionally equal to the abstract
semantics.  The engineering face --- how to shrink the feasible family
while preserving the queries that matter, and where compression
provably cannot go --- is owned by the shrinking note
(\code{ste-feasible-set-shrinking.tex}).  The database lineage the
paper builds on is conditional tables and provenance
\citep{imielinski1984incomplete,green2007provenance}, with
truth-maintenance ancestry \citep{doyle1979tms,dekleer1986atms}.

One theorem from that development must be quoted here, because
Chapter~\ref{ch:coref} is built directly on it.  Must-coreference is
the lattice meet of the feasible hypotheses' partitions:
\[\code{mustSetoid}\,(D) \;=\; \operatorname{sInf}\,
  \{\,\text{exact partition of } h : h \in \text{feasible}(D)\,\}
\]
(\code{mustSetoid\_eq\_sInf}, \Ste{PartitionRank}), monotone in the
corpus (\code{mustSetoid\_mono}, \code{rank\_mustSetoid\_mono}).  This
is Shannon's lattice reading of the STE consensus: what \emph{every}
surviving hypothesis agrees on.  Note its position in the spine: it is
a \emph{meet} over hypotheses --- an agreement question --- and the
dichotomy warns that such questions are the ill-conditioned ones.  The
frontier chapter will take up the dual, join-side notion of identity
and find it, by contrast, benign.

\chapter{Three corpora, three shapes}
\label{ch:empirical}

The theory has been run, end to end, on real text three times, with a
fallible LLM extractor producing set-valued frames and the verified
Lean pipeline doing everything from the frames onward.  The three runs
returned three different shapes of verdict, and together they span the
gluing trichotomy of Chapter~\ref{ch:gluing}.

\emph{The tides corpus} (multi-voice historical writing on tidal
theory \citep{galileo1632dialogue,kepler1609astronomia,
newton1687principia,descartes1644principia,pliny77naturalis}): a
consistent core exists after the machine \emph{corrects} extraction
slips --- the inconsistency detector doing its founding job of invalid
information detection.  \emph{The Federalist corpus}
\citep{hamilton1788federalist}: 86 documents pass through the proven
pipeline and cohere at scale --- the linear per-variable scan, proved
complete for rectangular frame corpora, verifying feasibility cheaply
exactly as the Helly-2 theory predicts.  \emph{The ratification
two-camp corpus}: the consistent core \emph{shatters}, cleanly and
camp-wise --- global gluing fails while camp-local coherence survives,
the gluing complex exhibited in the wild.  \textbf{The full treatments
are the three results notes}
(\code{ste-corpus-tides-results.tex},
\code{ste-corpus-federalist-results.tex},
\code{ste-corpus-ratification-twocamp.tex}), each of which states its
own boundary-of-trust section; the fidelity/consistency conditional ---
the kernel verifies coherence of the extracted frames, never fidelity
of extraction --- is stated there and in the epistemology note, and we
do not re-litigate it.

For the arc, the empirical chapters certify one thing: the non-metric
theory is not a redescription but a running instrument, and its three
possible verdicts --- cohere, cohere-after-correction, shatter-with-
located-obstruction --- are exactly the three outcomes the structure
theory said were possible.

\part{Language III: the Frontier}
\label{part:frontier}

\chapter{Multi-document coreference: identity as a lattice limit}
\label{ch:coref}

This chapter is the monograph's frontier and is originated here in
full: its mathematics exists only in the Lean module \Ste{Coreference}
and in the interpretive log, and this is its first complete prose
treatment.

\section{The problem, and why it is the right frontier}

Everything in Part~\ref{part:frames} lived, formally, inside one
mention space: the dynamic-frames development partitions a single
\code{Claim} type, and its must-coreference is a single partition of
it.  But the questions that motivate corpus-scale semantics are
cross-document by nature: \emph{is the ``General Carlson'' of document
12 the ``B.~Carlson'' of document 40?}  Cross-document coreference
resolution is a mature engineering problem
\citep{singh2011crossDocumentCoreference,barhom2021eventCoreference,
cattan2021sequentialCoreference}, entangled with the hard ontology of
event identity \citep{hovy2013events}, and closely allied to record
linkage and truth discovery in databases
\citep{elmagarmid2007duplicateRecordDetection,dong2009dataFusion,
li2015truthDiscoverySurvey,yin2007truthfinder}.  The dominant method
shape is the same everywhere: featurize pairs of mentions, score
similarity, cluster --- identity as the output of an optimizer.

The STE stance dictates a different decomposition of the problem, and
it is the same decomposition the whole book has practiced: separate
\emph{what the evidence forces} from \emph{what a scorer guesses}.  A
similarity model may propose links; that is extraction, fallible and
outside the kernel.  What the kernel can own is the question:
\emph{given} the per-document coreference decisions and a set of
asserted cross-document links, what identifications follow?  The answer
is not a cluster assignment but a partition-lattice construction ---
and it turns out to have a universal property, exact monotonicity laws,
a locality guardrail, and a place on the benign side of the spine's
dichotomy.  Each of these is a machine-checked theorem.

\section{Mentions with provenance}

Fix a type $\mathrm{Doc}$ of documents and, for each $d$, a type
$\mathrm{Local}\,d$ of mention identifiers local to $d$.  A
\emph{mention} is a dependent pair:
\[
  \mathrm{Mention} \;=\; \Sigma_{d : \mathrm{Doc}}\, \mathrm{Local}\,d
  \qquad (\code{Mention},\ \Ste{Coreference}).
\]
Provenance is thus carried in the type: two mentions from different
documents are never conflated by construction, only ever by
\emph{evidence}.  Each document contributes its local coreference
decisions as a setoid (equivalence relation)
$\mathrm{loc}\,d$ on $\mathrm{Local}\,d$ --- the output of
within-document resolution, which we take as given.  The embedding
\code{docSetoid} lifts $\mathrm{loc}\,d$ to a setoid on all of
$\mathrm{Mention}$ that relates mentions of document $d$ as
$\mathrm{loc}\,d$ does and relates every other mention only to itself.
The lift is provably provenance-respecting: it lies below the
same-document kernel setoid
$\code{sameDocSetoid} = \operatorname{ker}(\pi_1)$
(\code{docSetoid\_le\_sameDoc}).

Cross-document evidence enters as a \emph{raw relation}
$\mathrm{Link} \subseteq \mathrm{Mention}\times\mathrm{Mention}$ ---
deliberately not assumed symmetric, transitive, or reflexive.  This
models the epistemic situation honestly: an extractor asserts
individual links (``these two mentions co-refer''), and nothing about
the assertion process guarantees closure properties.  Equivalence is
something we \emph{impose} on evidence by closure, never something we
\emph{presume} of it.

\section{The forced coreference and its universal property}

\begin{definition}[\code{forcedCoref}, \Ste{Coreference}]
Given local coreferences $\mathrm{loc}$, links $\mathrm{Link}$, and an
active document set $S \subseteq \mathrm{Doc}$, the
\emph{corpus-forced coreference} is the join, in the complete lattice
of setoids on $\mathrm{Mention}$, of every active document's
contribution with the equivalence closure of the links:
\[\mathrm{forcedCoref}(\mathrm{loc},\mathrm{Link},S)
  \;=\;
  \bigsqcup_{d\in S} \mathrm{docSetoid}(\mathrm{loc},d)
  \;\sqcup\;
  \mathrm{EqvGen}(\mathrm{Link})
\]
\end{definition}

Call a setoid $s$ \emph{adequate} (to the evidence) if it dominates
every active document's contribution and contains every asserted link.
The construction has the universal property one should demand of
``what the evidence forces'':

\begin{theorem}[Identity as a lattice limit;
\code{forcedCoref\_eq\_sInf}]
The forced coreference is the least adequate setoid:
\[\mathrm{forcedCoref}(\mathrm{loc},\mathrm{Link},S)
  \;=\;
  \operatorname{sInf}\,\bigl\{\, s \;\bigm|\;
    \forall d\in S,\ \mathrm{docSetoid}(\mathrm{loc},d) \le s
    \ \wedge\ \mathrm{Link} \subseteq s \,\bigr\}
\]
\end{theorem}

\begin{interpretation}
``Who is the same as whom, across the corpus'' is not a search problem
whose answer must be guessed; it is a \emph{fixed point of constraint
accumulation} --- the unique smallest partition consistent with all the
local readings and all the asserted links.  The join reading (close up
the evidence) and the meet reading (intersect everything adequate)
provably coincide, so the construction is canonical: any resolver that
returns fewer identifications has ignored evidence, and any that
returns more has invented some.  Identity is \emph{derived, not
posited}.
\end{interpretation}

The theorem is the exact dual of the single-document consensus theorem
quoted in Chapter~\ref{ch:dynamic}: there, \code{mustSetoid\_eq\_sInf}
computes the \emph{meet over hypotheses} --- what every surviving world
agrees on; here, \code{forcedCoref\_eq\_sInf} computes the \emph{join
over evidence} --- the least commitment adequate to what was actually
asserted.  Both are $\operatorname{sInf}$ characterizations in the same
lattice $\mathrm{Setoid}$, indexed differently: by feasible worlds
there, by document provenance here.  This duality --- consensus as
meet, accumulation as join --- is the spine's dichotomy expressed
inside a single application, and we return to it in
\S\ref{sec:coref-spine}.

\section{Monotonicity: evidence only sharpens}

\begin{theorem}[\code{forcedCoref\_mono\_docs},
\code{forcedCoref\_mono\_link}]
The forced coreference is monotone in the active document set and in
the link relation: $S \subseteq S'$ implies
$\mathrm{forcedCoref}(S) \le \mathrm{forcedCoref}(S')$, and pointwise
$\mathrm{Link} \subseteq \mathrm{Link}'$ implies
$\mathrm{forcedCoref}(\mathrm{Link}) \le
 \mathrm{forcedCoref}(\mathrm{Link}')$.
\end{theorem}

\begin{theorem}[\code{rank\_forcedCoref\_mono\_docs}]
On finite corpora the monotonicity is quantitative: adding documents
can only raise the partition rank of the forced coreference,
\emph{i.e.}\ the number of merges the evidence has bought.
\end{theorem}

\begin{interpretation}
Identity can only sharpen as documents arrive.  This is the
cross-document analogue of information monotonicity in classical STE
--- more constraints, smaller feasible set --- but note the polarity
flip that provenance indexing produces: in the dynamic-frames world,
document insertion is \emph{antitone} on the feasible hypothesis set;
here it is \emph{monotone} on forced identifications.  Both statements
are the same phenomenon (evidence accumulates) read on opposite sides
of a Galois-style correspondence between hypotheses and commitments.
Retraction --- a document deleted, a link withdrawn --- is deliberately
outside this chapter's formalism: the dynamic-frames paper owns
deletion semantics, and grafting its deletion-expansion laws onto the
provenance-indexed setting is future work
(\S\ref{sec:coref-open}).
\end{interpretation}

\section{The locality guardrail}

A resolver that can merge anything must be trusted not to merge
everything.  The formal system earns that trust with a separation
theorem:

\begin{theorem}[\code{forcedCoref\_no\_links\_same\_doc}]
With no cross-document links asserted, the forced coreference never
crosses a document boundary: if
$\mathrm{forcedCoref}(\mathrm{loc}, \varnothing, S)$ relates $x$ and
$y$, then $x$ and $y$ belong to the same document.  (Via
\code{forcedCoref\_no\_links\_le\_sameDoc}: the whole relation lies
below the same-document kernel.)
\end{theorem}

\begin{interpretation}
The formalism refuses to invent cross-document identity out of thin
air.  However rich each document's internal coreference, and however
many documents are active, no identification leaks across a boundary
without a chain of asserted evidence.  This is the guardrail that
similarity-clustering architectures lack by construction: there,
any two sufficiently similar mentions may be merged, and the burden of
\emph{not} merging namesakes falls on the score.  Here the default is
separation; merging is what must be paid for, link by link.  In the
epistemology note's terms, the guardrail is part of the coherence
discipline: the kernel polices what follows from evidence, and the
fallibility of the links themselves is exactly where the
fidelity/consistency conditional applies.
\end{interpretation}

\section{The witness: transitivity is the inference engine}

The guardrail bounds the formalism from below; a concrete witness
bounds it from above, showing the closure genuinely infers.

\begin{theorem}[\code{witness\_cross\_doc\_forced}]
Take three documents with one mention apiece, no internal coreference,
and a chain of two links: document~0 to document~1, and document~1 to
document~2 --- but, provably, no direct 0--2 link
(\code{chainLink\_not\_direct}), and the endpoints in genuinely
different documents (\code{m0\_ne\_m2\_doc}).  The forced coreference
identifies the document-0 mention with the document-2 mention.
\end{theorem}

\begin{interpretation}
Together the guardrail and the witness bracket the system's inferential
power exactly: \emph{nothing} without a chain of evidence,
\emph{everything} a chain licenses.  All the inferential work is done
by transitive closure over asserted links --- which is also a precise
statement of risk.  One wrong link does not stay local: closure can
propagate it into a cascade of merges.  The monotonicity theorems then
say the cascade is at least \emph{orderly} (adding evidence never
un-merges), and the rank corollary meters its size.  A production
resolver built on this kernel would therefore treat link assertion as
the audited, expensive act --- precisely the inversion of the
similarity-clustering economy, where merging is cheap and auditing is
impossible.
\end{interpretation}

\section{The spine, one last time}
\label{sec:coref-spine}

The forced coreference is a \emph{join-side} construction: it
accumulates evidence by $\sqcup$ in the setoid lattice.  The spine's
contraction theorem (\code{shannonDist\_sup\_right\_le},
Chapter~\ref{ch:spine}) holds in every setoid lattice over a finite
carrier, so it applies verbatim to the mention partitions of a finite
corpus: perturbing one contribution --- swapping a document's local
coreference for a nearby one, adding or dropping a link --- moves the
forced coreference by no more than the perturbation, in the
combinatorial Shannon metric.  Cross-document identity, as formalized
here, is a \emph{well-conditioned} quantity.

The meet-side questions are also present in this domain, and the
dichotomy predicts their character.  ``What identifications do
\emph{all} feasible hypotheses agree on'' (the \code{mustSetoid}
consensus of Chapter~\ref{ch:dynamic}) and ``what do these two
corpora's coreference structures have in common'' are
$\sqcap$-questions, and \code{shannonDist\_inf\_not\_contraction} warns
that they are ill-conditioned: a small change in one input can move the
agreement arbitrarily far.  The practical reading is a design rule for
corpus systems: \emph{publish forced (join-side) identity; treat
consensus (meet-side) identity as fragile, and never build downstream
pipelines that differentiate through it}.  That rule is this
monograph's thesis made operational, and it is the reason the frontier
chapter belongs to the same book as the geometry chapter.

\section{What is deliberately not yet mechanized}
\label{sec:coref-open}

The chapter closes with its honest boundary, in the project's standard
form.  Mechanized: everything cited above, sorry-free.  Not mechanized,
and flagged as the frontier's frontier:
(i)~\emph{link fallibility} --- graded links with a fuzzy-frame lift
(Chapter~\ref{ch:frames}), so that a forced merge carries the minimum
confidence along its evidence chain, in the spirit of subjective logic
\citep{josang2016subjectiveLogic};
(ii)~\emph{retraction} --- transporting the dynamic-frames deletion
laws to the provenance-indexed setting;
(iii)~\emph{entity-level frames} --- attaching set-valued frames to
merged clusters and asking when the merged frame remains feasible,
which is where coreference re-enters the gluing theory of
Chapter~\ref{ch:gluing};
(iv)~\emph{the Helly question} --- whether 2-locality survives the
sigma-type indexing of mentions, or whether provenance, like
second-order possibility, has a locality price.  Conjecture: it
survives for link-free corpora (which are rectangular in the
appropriate sense) and fails in general once links couple documents ---
that is, the dichotomy will reappear exactly here.

\chapter{The honest boundary}
\label{ch:boundary}

Every load-bearing caveat in this book has a single owning statement in
the epistemology note (\code{hyperframes-epistemology.tex}), and this
chapter is only its index.  The theory studies \emph{coherence, not
truth}; its verdicts are conditional on extraction fidelity, and the
kernel checks the consistency side of that conditional only.  The crisp
set-valued formalism has a stated expressiveness ceiling --- no
normalized densities, hence no genuine Shannon entropy over readings,
until the mechanization landscape changes; the fuzzy stack of
Chapter~\ref{ch:frames} is the possibilistic approximation, honestly
labeled.  The empirical chapters' boundary-of-trust sections
instantiate the conditional per corpus.  Nothing in this monograph
strengthens, weakens, or renegotiates that position; a reader who wants
it argued should read the note, which argues it at length.

\appendix

\chapter{The overlap-resolution map}
\label{app:map}

This appendix is the monograph's contract with its corpus: for every
topic, exactly one document owns the deep dive; this book cites owners
and never re-derives; and the residual overlaps among existing
documents are flagged for trimming, as recommendations only --- no
existing paper is edited by this book.

\section{Ownership}

\begin{center}
\small
\begin{tabular}{@{}p{0.44\columnwidth}p{0.5\columnwidth}@{}}
\toprule
\textbf{Topic} & \textbf{Owning deep dive} \\
\midrule
Formal STE core; taxonomy; Carlson counting/reduction chain &
  \code{ste-mechanization.tex} \\
Primary-source review (Combettes, Carlson) &
  \code{literature-review.tex} \\
Projection/convex inheritance; set-membership tradition &
  \code{topological-ste-survey.tex} \\
Post-2012 uptake; mechanization gap &
  \code{decade-followon-survey.tex} \\
Sheaf/\v{C}ech gluing for STE; support--cover theorem &
  \code{ste-cohomology.tex} (flagship) \\
Accessible gluing dictionary &
  \code{ste-cohomology-recast.tex} \\
Contextuality-literature setting of the obstruction &
  \code{cohomological-obstruction-litreview.tex} \\
Representation, treewidth stack, exponential lower bounds &
  \code{ste-representation-sheaves.tex} \\
Obstruction-vs-blow-up sandwich, literature &
  \code{representation-bounds-litreview.tex} \\
CG-as-STE equivalence and its scope &
  \code{constraint-grammar-ste.tex} \\
Historical bridge to CL; FCA/hyperframe lattice &
  \code{ste-and-language.tex} \\
Epistemology; coherence-not-truth; entropy ceiling &
  \code{hyperframes-epistemology.tex} \\
Dynamic corpus semantics; modal identity grades; $2^N$ limit &
  \code{dynamic-frame-normalization.tex} \\
Feasible-family reduction/compression contract &
  \code{ste-feasible-set-shrinking.tex} \\
New Lean threads (set-valued, fuzzy, set-of-sets, gluing complex,
info-geometry): interpretations &
  \code{hyperframe-interpretations.tex} (running log) \\
Tides / Federalist / ratification runs &
  the three \code{ste-corpus-*-results} notes \\
Multi-document coreference (\Ste{Coreference}) &
  \textbf{this monograph}, Chapter~\ref{ch:coref} \\
Conditioning dichotomy as unifying thread &
  \textbf{this monograph}, Chapter~\ref{ch:spine} \\
\bottomrule
\end{tabular}
\end{center}

\section{Overlaps flagged for trimming (recommendations)}

\begin{enumerate}
\item \textbf{Flagship vs.\ representation paper.}
  \code{ste-cohomology.tex} currently re-develops the
  bounded-treewidth algorithmic stack and the rectangle-cover /
  fooling-set bounds that \code{ste-representation-sheaves.tex} owns.
  Recommend: the flagship trims those sections to statements plus
  citations, keeping only what the support--cover theorem itself needs.
\item \textbf{Flagship vs.\ mechanization paper.}
  The flagship's foundations sections (\S2--3) re-explain the STE core
  and the Carlson turn owned by \code{ste-mechanization.tex} and the
  surveys.  Recommend: compress to a one-page recap with pointers.
\item \textbf{Recast note vs.\ the metric-geometry thread.}
  \code{ste-cohomology-recast.tex}'s ``readings form a metric space''
  section predates \Ste{InfoTopology} and now duplicates (in weaker
  form) material owned by the interpretive log and
  Chapter~\ref{ch:spine}.  Recommend: reduce to a pointer.
\item \textbf{The three surveys.}
  \code{literature-review.tex}, \code{topological-ste-survey.tex}, and
  \code{decade-followon-survey.tex} each re-summarize
  \citet{carlson2012}'s core construction.  Recommend: the two surveys
  cite the literature review's summary instead of repeating it, keeping
  only their period-specific material.
\item \textbf{Dynamic frames vs.\ shrinking note.}
  The $2^N$ tight-limit theorem, minimal supports, and cross-corpus
  gluing appear in both \code{dynamic-frame-normalization.tex} and
  \code{ste-feasible-set-shrinking.tex}.  Recommend: the shrinking note
  keeps the practitioner contract and cites the paper for the theorems.
\item \textbf{CG note vs.\ language note.}
  \code{constraint-grammar-ste.tex}'s ``consequences transported''
  section partially re-derives monotonicity facts owned by
  \code{ste-and-language.tex}.  Minor; recommend cross-reference.
\item \textbf{Tides plan vs.\ tides results.}
  \code{ste-corpus-tides-plan.tex} is superseded by the results note.
  Recommend: mark as historical/superseded in its header.
\item \textbf{Gluing complex placement.}
  \Ste{LocalGluing} currently has prose only in the log and in
  Chapter~\ref{ch:gluing} here.  Recommend: its deep dive be added to a
  future revision of the flagship (not to a new paper), keeping
  ownership consolidated.
\end{enumerate}

\chapter{Index of machine-checked results cited}
\label{app:lean}

All identifiers verified sorry-free in the project's Lean~4 library at
the time of writing; module names are relative to \code{lean/Ste/}.

\begin{center}
\footnotesize
\begin{tabular}{@{}>{\raggedright\arraybackslash}p{0.34\columnwidth}>{\raggedright\arraybackslash}p{0.28\columnwidth}>{\raggedright\arraybackslash}p{0.24\columnwidth}@{}}
\toprule
\textbf{Lean identifier} & \textbf{Module} & \textbf{Cited in} \\
\midrule
\code{frame\_hasHelly\_two} & \Ste{HellyConsistency} & Ch.~\ref{ch:spine} \\
\code{rank\_submodular} & \Ste{PartitionRank} & Ch.~\ref{ch:spine} \\
\code{mustSetoid\_eq\_sInf} & \Ste{PartitionRank} & Ch.~\ref{ch:dynamic} \\
\code{mustSetoid\_mono} & \Ste{PartitionRank} & Ch.~\ref{ch:dynamic} \\
\code{rank\_mustSetoid\_mono} & \Ste{PartitionRank} & Ch.~\ref{ch:dynamic} \\
\code{couplingExcess\_pos\_of\_not\_le} & \Ste{CouplingRank} & Ch.~\ref{ch:spine} \\
\code{shannonDist}, \code{shannonDist\_triangle} & \Ste{InfoDistance} & Ch.~\ref{ch:spine} \\
\code{one\_le\_shannonDist\_of\_ne} & \Ste{InfoTopology} & Ch.~\ref{ch:spine} \\
\code{shannonDist\_sup\_right\_le} & \Ste{InfoTopology} & Ch.~\ref{ch:spine}, \ref{ch:coref} \\
\code{shannonDist\_sup\_le\_add} & \Ste{InfoTopology} & Ch.~\ref{ch:spine} \\
\code{shannonDist\_inf\_not\_contraction} & \Ste{InfoTopology} & Ch.~\ref{ch:spine}, \ref{ch:coref} \\
\code{shannonDist\_add\_rank\_add\_rank} & \Ste{InfoTopology} & Ch.~\ref{ch:spine} \\
\code{pentagonal\_violation}, \code{median\_not\_unique} & \Ste{InfoTopology} & Ch.~\ref{ch:spine} \\
\code{aligned\_of\_chain}, \code{le\_sup\_of\_aligned} & \Ste{InfoTopology} & Ch.~\ref{ch:spine} \\
\code{diagonal\_cechObstruction} & \Ste{CechObstruction} & Ch.~\ref{ch:spine}, \ref{ch:gluing} \\
\code{GluesOn}, \code{gluingComplex} & \Ste{LocalGluing} & Ch.~\ref{ch:gluing} \\
\code{mem\_gluingComplex\_of\_subset} & \Ste{LocalGluing} & Ch.~\ref{ch:gluing} \\
\code{cechVanishesCover\_iff\_gluesLocallyAt\_card} & \Ste{LocalGluing} & Ch.~\ref{ch:gluing} \\
\code{diagonal\_gluesLocallyAt\_one} & \Ste{LocalGluing} & Ch.~\ref{ch:gluing} \\
\code{diagonal\_not\_gluesLocallyAt\_two} & \Ste{LocalGluing} & Ch.~\ref{ch:gluing} \\
\code{diagonal\_local\_not\_tight} & \Ste{LocalGluing} & Ch.~\ref{ch:gluing} \\
\code{cechVanishesCover\_of\_hasHelly\_gluingFiber} & \Ste{LocalGluing} & Ch.~\ref{ch:gluing} \\
\code{propertySet\_isRectangle} & \Ste{SetValuedFrame} & Ch.~\ref{ch:frames} \\
\code{propertySet\_cechVanishesCover} & \Ste{SetValuedFrame} & Ch.~\ref{ch:spine}, \ref{ch:frames} \\
\code{corpusFeasibilitySet\_nonempty\_iff} & \Ste{SetValuedFrame} & Ch.~\ref{ch:frames} \\
\code{isExtent\_corpusFeasibilitySet} & \Ste{SetValuedFrame} & Ch.~\ref{ch:frames} \\
\code{hasHelly\_ofCrisp\_two} & \Ste{SetValuedFrame} & Ch.~\ref{ch:frames} \\
\code{consistencyDegree} & \Ste{FuzzyFrame} & Ch.~\ref{ch:frames} \\
\code{cut\_fuzzyFeasibility} & \Ste{FuzzyFrame} & Ch.~\ref{ch:frames} \\
\code{consistencyDegree\_crisp\_eq\_one\_iff} & \Ste{FuzzyFrame} & Ch.~\ref{ch:frames} \\
\code{weightedFrame\_gradedHelly\_two} & \Ste{FuzzyFrame} & Ch.~\ref{ch:frames} \\
\code{SelectableConsistent} & \Ste{Carlson.SetOfSets} & Ch.~\ref{ch:frames} \\
\code{selectable\_singleton\_iff} & \Ste{Carlson.SetOfSets} & Ch.~\ref{ch:frames} \\
\code{hasHellySel\_singletonFamily\_iff} & \Ste{Carlson.SetOfSets} & Ch.~\ref{ch:frames} \\
\code{exists\_not\_hasHellySel\_two} & \Ste{Carlson.SetOfSets} & Ch.~\ref{ch:spine}, \ref{ch:frames} \\
\code{Mention}, \code{docSetoid}, \code{sameDocSetoid} & \Ste{Coreference} & Ch.~\ref{ch:coref} \\
\code{docSetoid\_le\_sameDoc} & \Ste{Coreference} & Ch.~\ref{ch:coref} \\
\code{forcedCoref} & \Ste{Coreference} & Ch.~\ref{ch:coref} \\
\code{forcedCoref\_eq\_sInf} & \Ste{Coreference} & Ch.~\ref{ch:coref} \\
\code{forcedCoref\_mono\_docs} & \Ste{Coreference} & Ch.~\ref{ch:coref} \\
\code{forcedCoref\_mono\_link} & \Ste{Coreference} & Ch.~\ref{ch:coref} \\
\code{forcedCoref\_no\_links\_le\_sameDoc} & \Ste{Coreference} & Ch.~\ref{ch:coref} \\
\code{forcedCoref\_no\_links\_same\_doc} & \Ste{Coreference} & Ch.~\ref{ch:coref} \\
\code{witness\_cross\_doc\_forced} & \Ste{Coreference} & Ch.~\ref{ch:coref} \\
\code{chainLink\_not\_direct}, \code{m0\_ne\_m2\_doc} & \Ste{Coreference} & Ch.~\ref{ch:coref} \\
\code{rank\_forcedCoref\_mono\_docs} & \Ste{Coreference} & Ch.~\ref{ch:coref} \\
\bottomrule
\end{tabular}
\end{center}

\bibliographystyle{plainnat}
\bibliography{../../refs}

\end{document}
